\VerdierPhysicalPageAnchor{0054}%
\VerdierSemanticAnchor{ch1-1-5-4-eq1-condition}
soient $\alpha(\lambda,\mu)$-commutatifs lorsque
$\lambda,\mu\in\mathbb Z^{n-1}$ ou
$\lambda,\mu\in\mathbb Z e_n$. Désignons alors par
$Q(\sigma,\tau)$, $\sigma\in\mathbb Z^{n-1}$,
$\tau\in\mathbb Z e_n$, la propriété :

\VerdierSemanticAnchor{ch1-1-5-4-property-q}
\begin{quote}
  \noindent\textup{$(Q(\sigma,\tau))$}\quad
  \emph{Le diagramme ci-après est
  $\alpha(\sigma,\tau)\alpha^{-1}(\tau,\sigma)$-commutatif :}
\end{quote}

\VerdierSemanticAnchor{ch1-1-5-4-eq2}
\begin{equation}
  \resizebox{0.98\textwidth}{!}{%
    $\begin{tikzcd}[ampersand replacement=\&,column sep=4.5em,row sep=3.2em]
      FT(\sigma+\tau)
        \arrow[r,"{F*c(\sigma,\tau)}"]
        \arrow[d,"{F*c(\tau,\sigma)}"]
        \& FT(\tau)T(\sigma)
        \arrow[r,"{p(\tau)*T(\sigma)}"]
        \& T'(\tau)FT(\sigma)
        \arrow[d,"{T'(\tau)*p(\sigma)}"'] \\
      FT(\sigma)T(\tau)
        \arrow[d,"{p(\sigma)*T(\tau)}"]
        \& {}
        \& T'(\tau)T'(\sigma)F
        \arrow[d,"{c'(\sigma,\tau)^{-1}*F}"'] \\
      T'(\sigma)FT(\tau)
        \arrow[r,"{T'(\sigma)*p(\tau)}"']
        \& T'(\sigma)T'(\tau)F
        \arrow[r,"{c'(\tau,\sigma)^{-1}*F}"']
        \& T'(\sigma+\tau)F
    \end{tikzcd}$%
  }\;.
  \tag*{\(\textup{(1.5.4.2)}_{\sigma,\tau}\)}
\end{equation}

Par hypothèse, les propriétés
$Q(e_i,e_n)$, $Q(e_i,0)$, $1\leq i\leq n-1$, $Q(0,e_n)$,
sont vraies. De plus, en utilisant les diagrammes
$\alpha(\lambda,\mu)$-commutatifs
$\textup{(1.5.4.1)}_{\lambda,\mu}$, on vérifie aisément
(1.5.4.3) :

\VerdierSemanticAnchor{ch1-1-5-4-eq3}
\begin{quote}
  \noindent\textup{(1.5.4.3)}\quad
  \emph{Si deux des propriétés
  $Q(\sigma,\tau+\tau')$, $Q(\sigma,\tau)$, $Q(\sigma,\tau')$,
  $\sigma\in\mathbb Z^{n-1}$, $\tau,\tau'\in\mathbb Z e_n$
  sont vérifiées, la troisième l'est aussi. De même, si deux des
  propriétés $Q(\sigma+\sigma',\tau)$, $Q(\sigma,\tau)$,
  $Q(\sigma',\tau)$, $\sigma,\sigma'\in\mathbb Z^{n-1}$,
  $\tau\in\mathbb Z e_n$ sont vraies, la troisième l'est aussi.}
\end{quote}

\VerdierSemanticAnchor{ch1-1-5-4-q-conclusion}
On en déduit que la propriété $Q(\sigma,\tau)$ est vérifiée pour tout
$\sigma\in\mathbb Z^{n-1}$ et tout $\tau\in\mathbb Z e_n$.

\VerdierSemanticAnchor{ch1-1-5-4-p-omega}
Définissons alors, pour tout $\omega\in\mathbb Z^n$,
$\omega=\sigma+\tau$, $\sigma\in\mathbb Z^{n-1}$,
$\tau\in\mathbb Z e_n$, l'isomorphisme
$p(\omega)\colon FT(\omega)\longrightarrow T'(\omega)F$, comme étant
l'unique isomorphisme rendant le diagramme
$\textup{(1.5.4.1)}_{\sigma,\tau}$
$\alpha(\sigma,\tau)$-commutatif. Ceci nous permet d'écrire pour tout
couple $\lambda,\mu\in\mathbb Z^n$ le diagramme
$\textup{(1.5.4.1)}_{\lambda,\mu}$ et de considérer la propriété :

\VerdierSemanticAnchor{ch1-1-5-4-property-p}
\begin{quote}
  \noindent\textup{$(P(\lambda,\mu))$}\quad
  \emph{Le diagramme $\textup{(1.5.4.1)}_{\lambda,\mu}$ est
  $\alpha(\lambda,\mu)$-commutatif.}
\end{quote}
